\documentclass[11pt]{article}

\usepackage{fullpage}
\usepackage{amsmath}
\usepackage{amssymb}
\usepackage{amsfonts}
\usepackage{amsthm}
\usepackage{mathtools}
\usepackage{enumitem}
\usepackage[colorlinks,linkcolor=magenta,citecolor=blue]{hyperref}

\newcommand{\Rbar}{\bar R}
\newcommand{\Sbar}{\bar S}
\newcommand{\Om}{\Omega}
\newcommand{\E}{\mathbb E}
\newcommand{\Var}{\operatorname{Var}}
\newcommand{\tr}{\operatorname{tr}}
\newcommand{\diag}{\operatorname{diag}}
\newcommand{\Unif}{\operatorname{Unif}}
\newcommand{\Gam}{\operatorname{Gamma}}
\newcommand{\Wishart}{\mathcal W}
\newcommand{\Normal}{\mathcal N}

\newtheorem{lemma}{Lemma}

\title{A Wishart-column alternative to the inverse-Wishart LD relaxation}
\author{}
\date{}

\begin{document}
\maketitle

\section{Setup}

The notes in \texttt{susie-relax-report.tex} and
\texttt{summary-susie-love-revisit.tex} use an inverse-Wishart prior on
the latent LD matrix. This note considers the corresponding alternative
\[
R\sim \Wishart_J\left(\frac{\Rbar}{\nu_0},\nu_0\right),
\]
so that \(\E(R)=\Rbar\). Here \(\Wishart(V,\nu)\) denotes the Wishart
distribution with scale matrix \(V\) and degrees of freedom \(\nu\).
For a nonsingular \(J\times J\) Wishart distribution, we need
\(\nu_0>J-1\). This is a practical difference from the inverse-Wishart
parameterization, where \(\nu_0\) can be used more directly as a small
positive relaxation parameter.

Throughout, \(x\in\mathbb R^J\) is the vector of marginal summary
statistics, \(N\) is the GWAS sample size, and \(\Rbar\) is a reference
LD correlation matrix. For a fixed coordinate \(j\), write
\[
\mu_j=\Rbar_{-j,j},\qquad
\Sbar_j=\Rbar_{-j,-j}-\mu_j\mu_j^\top,\qquad
\Omega=\Rbar^{-1}.
\]
Let
\[
d_j=R_{jj},\qquad
z_j=\frac{R_{-j,j}}{R_{jj}},\qquad
S_j^R=R_{-j,-j}-R_{-j,j}R_{jj}^{-1}R_{j,-j}.
\]
The active column of \(R\) is \(d_j u_j(z_j)\), where \(u_j(z_j)\) is the
vector with \(j\)-th entry equal to \(1\) and remaining entries equal to
\(z_j\).

\section{Wishart-induced column distribution}

By the standard block decomposition of a Wishart random matrix, after
permuting coordinate \(j\) to the first position,
\begin{align}
d_j
&\sim
\Gam\left(\frac{\nu_0}{2},\frac{\nu_0}{2}\right), \\
z_j\mid d_j
&\sim
\Normal_{J-1}\left(
\mu_j,\frac{\Sbar_j}{\nu_0 d_j}
\right), \\
S_j^R
&\sim
\Wishart_{J-1}\left(\frac{\Sbar_j}{\nu_0},\nu_0-1\right),
\end{align}
using the shape--rate parameterization for Gamma distributions. Moreover,
\(S_j^R\) is independent of \((d_j,z_j)\). These facts imply
\[
\E(d_j)=1,\qquad
\Var(d_j)=\frac{2}{\nu_0},\qquad
\E(z_j\mid d_j)=\mu_j,\qquad
\Var(z_j\mid d_j)=\frac{\Sbar_j}{\nu_0 d_j}.
\]
The mean of the Schur complement is
\[
\E(S_j^R)=\frac{\nu_0-1}{\nu_0}\Sbar_j.
\]
This is not equal to \(\Sbar_j\), but the full matrix still has
\(\E(R)=\Rbar\), because
\[
\E\{d_j z_jz_j^\top\}
=
\mu_j\mu_j^\top+\frac{1}{\nu_0}\Sbar_j.
\]

\paragraph{Comparison with the inverse-Wishart column.}
Under the inverse-Wishart model used in the earlier notes, \(d_j\) is
inverse-gamma and the relaxed column direction has a normal scale tied to
the random Schur complement. Under the Wishart model, \(d_j\) is gamma
and
\[
z_j\mid d_j\sim
\Normal_{J-1}\left(\mu_j,\frac{\Sbar_j}{\nu_0d_j}\right),
\]
so the column direction is tied to the reference Schur complement
\(\Sbar_j\), not to \(S_j^R\). This changes the conjugacy structure.

\section{Fixed-variance SER-relax with a Wishart column}

The closest analogue of the prototype in
\texttt{susie-relax-report.tex} keeps the observation covariance fixed at
\(\Rbar/N\), but replaces the inverse-Wishart column prior by the
Wishart-induced column prior:
\begin{align}
x\mid \beta,\gamma=j,d_j,z_j
&\sim
\Normal_J\left(\beta d_j u_j(z_j),\frac{\Rbar}{N}\right), \\
\beta&\sim \Normal(0,\sigma_0^2), \\
\gamma&\sim \Unif\{1,\ldots,J\}, \\
d_j&\sim \Gam\left(\frac{\nu_0}{2},\frac{\nu_0}{2}\right), \\
z_j\mid d_j
&\sim
\Normal_{J-1}\left(\mu_j,\frac{\Sbar_j}{\nu_0d_j}\right).
\end{align}
This model is computationally close to the inverse-Wishart prototype, but
the factor \(q_j(d_j)\) changes because \(d_j\) appears in the likelihood
and in the precision of the \(z_j\) prior.

\paragraph{Integrating out \(z_1\).}
For concreteness, take \(\gamma=1\) and write
\[
x=
\begin{pmatrix}
x_1\\x_{-1}
\end{pmatrix},
\qquad
\Rbar=
\begin{pmatrix}
1 & \mu_1^\top\\
\mu_1 & \Sbar_1+\mu_1\mu_1^\top
\end{pmatrix}.
\]
Conditional on \((\beta,d_1)\), the only random part of the mean is
\(\beta d_1z_1\), and
\[
z_1\mid d_1
\sim
\Normal_{J-1}\left(\mu_1,\frac{\Sbar_1}{\nu_0d_1}\right).
\]
Thus, after integrating over \(z_1\),
\[
x\mid \beta,d_1,\gamma=1
\sim
\Normal_J\left(
\beta d_1
\begin{pmatrix}
1\\\mu_1
\end{pmatrix},
\frac{\Rbar}{N}
+\frac{\beta^2d_1}{\nu_0}
\begin{pmatrix}
0 & 0\\
0 & \Sbar_1
\end{pmatrix}
\right).
\]
Equivalently, using the block normal conditional,
\begin{align}
p(x\mid \beta,d_1,\gamma=1)
&=
\Normal\left(x_1\mid \beta d_1,\frac{1}{N}\right)
\nonumber\\
&\quad\times
\Normal_{J-1}\left(
x_{-1}\mid x_1\mu_1,\,
\left(\frac{1}{N}+\frac{\beta^2d_1}{\nu_0}\right)\Sbar_1
\right).
\end{align}
This likelihood keeps the same conditional-mean simplification:
after conditioning on \(x_1\), the mean of \(x_{-1}\) is \(x_1\mu_1\)
and does not involve \(\beta\) or \(d_1\).

Now write the same likelihood explicitly in terms of
\[
q=x^\top\Rbar^{-1}x,
\qquad
\zeta=\frac{N\beta^2d_1}{\nu_0}.
\]
The covariance matrix above is \(N^{-1}\widetilde R\), where
\[
\widetilde R
=
\Rbar+\zeta
\begin{pmatrix}
0 & 0\\
0 & \Sbar_1
\end{pmatrix}.
\]
By the block determinant identity,
\[
|\widetilde R|=(1+\zeta)^{J-1}|\Rbar|,
\]
and by the block inverse identity,
\[
x^\top\widetilde R^{-1}x
=
\frac{1}{1+\zeta}x^\top\Rbar^{-1}x
+\frac{\zeta}{1+\zeta}x_1^2.
\]
Also, since \((1,\mu_1^\top)^\top\) is the first column of
\(\widetilde R\),
\[
\widetilde R^{-1}
\begin{pmatrix}
1\\\mu_1
\end{pmatrix}
=e_1.
\]
Therefore
\begin{align}
p(x\mid \beta,d_1,\gamma=1)
&=
\frac{N^{J/2}}
{(2\pi)^{J/2}|\Rbar|^{1/2}(1+\zeta)^{(J-1)/2}}
\nonumber\\
&\quad\times
\exp\left\{
-\frac{N}{2}
\left[
\frac{q}{1+\zeta}
+\frac{\zeta}{1+\zeta}x_1^2
-2\beta d_1x_1
+\beta^2d_1^2
\right]
\right\}.
\end{align}
Equivalently,
\begin{align}
p(x\mid \beta,d_1,\gamma=1)
&\propto
(1+\zeta)^{-(J-1)/2}
\exp\left\{-\frac{Nq}{2(1+\zeta)}\right\}
\nonumber\\
&\quad\times
\exp\left\{-\frac{N\zeta x_1^2}{2(1+\zeta)}\right\}
\exp\left\{N\beta d_1x_1-\frac{N}{2}\beta^2d_1^2\right\},
\end{align}
where the omitted proportionality constant is common across coordinates.

\subsection{Mean-field updates}

Use
\[
q(\gamma=j,\beta,d_j,z_j)
=
\pi_jq_j(\beta)q_j(d_j)q_j(z_j),
\]
where
\[
q_j(\beta)=\Normal(m_{\beta j},v_{\beta j}),\qquad
q_j(z_j)=\Normal_{J-1}(m_{zj},V_{zj}).
\]
Let
\[
\bar\beta_j=m_{\beta j},\qquad
\overline{\beta^2}_j=v_{\beta j}+m_{\beta j}^2,\qquad
\bar d_j=\E_j(d_j),\qquad
\overline{d_j^2}=\E_j(d_j^2).
\]
Define
\begin{align}
A_j
&=
1+\tr(\Sbar_j^{-1}V_{zj})
+(m_{zj}-\mu_j)^\top \Sbar_j^{-1}(m_{zj}-\mu_j),\\
B_j
&=
x_j+(m_{zj}-\mu_j)^\top
\Sbar_j^{-1}(x_{-j}-\mu_jx_j),\\
Q_j
&=
\tr(\Sbar_j^{-1}V_{zj})
+(m_{zj}-\mu_j)^\top \Sbar_j^{-1}(m_{zj}-\mu_j).
\end{align}
The Gaussian update for the effect size is
\begin{align}
v_{\beta j}
&=
\left(\sigma_0^{-2}+N\overline{d_j^2}A_j\right)^{-1},\\
m_{\beta j}
&=
v_{\beta j}N\bar d_jB_j.
\end{align}
The Gaussian update for the column direction is
\begin{align}
V_{zj}
&=
\frac{\Sbar_j}
{N\overline{\beta^2}_j\overline{d_j^2}+\nu_0\bar d_j},\\
m_{zj}
&=
\mu_j+
\frac{N\bar\beta_j\bar d_j}
{N\overline{\beta^2}_j\overline{d_j^2}+\nu_0\bar d_j}
(x_{-j}-\mu_jx_j).
\end{align}
The optimal density for \(d_j\) is non-standard:
\[
q_j(d_j)
\propto
d_j^{\nu_0/2-1+(J-1)/2}
\exp\left\{
-\left(\frac{\nu_0}{2}+\frac{\nu_0 Q_j}{2}\right)d_j
-\frac{N}{2}\overline{\beta^2}_jA_jd_j^2
+N\bar\beta_jB_jd_j
\right\},
\qquad d_j>0.
\]
In terms of \(r_j=\log d_j\), including the Jacobian,
\[
\ell_j(r)
=
\left(\frac{\nu_0}{2}-1+\frac{J-1}{2}\right)r
-\left(\frac{\nu_0}{2}+\frac{\nu_0 Q_j}{2}\right)e^r
-\frac{N}{2}\overline{\beta^2}_jA_je^{2r}
+N\bar\beta_jB_je^r
+r.
\]
One-dimensional quadrature over \(r\) gives
\(\E_j(d_j)\), \(\E_j(d_j^2)\), \(\E_j(\log d_j)\), and
\(E_j\log q_j(d_j)\).

\subsection{Expected log terms}

The expected log-likelihood has the same form as in the inverse-Wishart
prototype:
\begin{align}
E_j\log p(x\mid \beta,d_j,z_j,\gamma=j)
&=
-\frac{J}{2}\log(2\pi)
-\frac12\log\left|\frac{\Rbar}{N}\right| \nonumber\\
&\quad
-\frac{N}{2}
\left[
x^\top\Omega x
-2\bar\beta_j\bar d_jB_j
+\overline{\beta^2}_j\overline{d_j^2}A_j
\right].
\end{align}
The Wishart-column prior contributions are
\begin{align}
E_j\log p(d_j\mid\nu_0)
&=
\frac{\nu_0}{2}\log\frac{\nu_0}{2}
-\log\Gamma\left(\frac{\nu_0}{2}\right)
+\left(\frac{\nu_0}{2}-1\right)\E_j\log d_j
-\frac{\nu_0}{2}\E_j d_j,\\
E_j\log p(z_j\mid d_j,\nu_0)
&=
-\frac{J-1}{2}\log(2\pi)
-\frac12\log|\Sbar_j|
+\frac{J-1}{2}\log\nu_0
+\frac{J-1}{2}\E_j\log d_j
-\frac{\nu_0}{2}\E_jd_j\,Q_j.
\end{align}
These terms can be substituted into the usual local lower bound and used
to update the coordinate probabilities \(\pi_j\).

\section{Takeaways}

\begin{enumerate}[leftmargin=*]
\item A Wishart prior \(R\sim\Wishart(\Rbar/\nu_0,\nu_0)\) induces
\[
d_j\sim\Gam(\nu_0/2,\nu_0/2),\qquad
z_j\mid d_j\sim
\Normal_{J-1}(\mu_j,\Sbar_j/(\nu_0d_j)).
\]
\item In the fixed-variance SER-relax analogue, variational inference is
nearly the same as the inverse-Wishart prototype, except \(q_j(d_j)\)
has gamma-type behavior near zero and the \(z_j\) prior contributes
\(\nu_0 d_j Q_j/2\).
\item After integrating out \(z_1\) in the fixed-variance model, the
likelihood factorizes as
\[
p(x\mid \beta,d_1,\gamma=1)
=
\Normal\left(x_1\mid \beta d_1,\frac{1}{N}\right)
\Normal_{J-1}\left(
x_{-1}\mid x_1\mu_1,\,
\left(\frac{1}{N}+\frac{\beta^2d_1}{\nu_0}\right)\Sbar_1
\right).
\]
\end{enumerate}

\end{document}
